% PAGES: 191-192
% Continues the OPEN "\begin{itemize}" left unclosed at the end of
% sec06_c2_3.tex (folio 174), whose last item ends "...cf. (2.19),". This
% file supplies the final item and "\end{itemize}" below on folio 175. Do
% not open a new "\begin{itemize}".

\item for a total operation count of
\[
\left(\frac{1}{3}n^3 + O(n^2)\right) \;+\; (n^2+O(n)) \;+\; (n^2+O(n))
\;=\; \frac{1}{3}n^3 + O(n^2);
\]
see (10.87) from Section 10.2.3 for a proof of the last equality.
\end{itemize}

\subsection{Solving linear systems corresponding to the same spd matrix}

In this section, we show how the Cholesky decomposition of a matrix can be used to
solve multiple linear systems corresponding to the same symmetric positive definite
matrix efficiently.

Assume that we want to solve $p$ linear systems corresponding to an $n \times n$
symmetric positive definite matrix $A$. In other words, we want to find $n \times 1$
vectors $x_i$, $i = 1:p$, such that
\begin{equation}
Ax_i \;=\; b_i, \quad \forall\, i = 1:p,
\end{equation}
where $b_i$ is a column vector of size $n$, for $i = 1:p$.

One way to solve the linear systems from (6.50) would be to use the routine
linear\_solve\_cholesky from Table 6.2 to solve each one of the $p$ linear systems
inside a ``for'' loop as follows:
\begin{center}
\fbox{\begin{minipage}{0.6\linewidth}
for $i = 1:p$\\
\hspace*{1.5em}$x_i = \text{linear\_solve\_cholesky}(A,b_i)$\\
end
\end{minipage}}
\end{center}

This would require
\begin{equation}
p\left(\frac{1}{3}n^3 + O(n^2)\right) \;=\; \frac{1}{3}pn^3 + pO(n^2)
\end{equation}
operations, since each linear solver requires $\frac{1}{3}n^3 + O(n^2)$ operations; cf.
Lemma 6.2.

However, the most expensive part of $x_i = \text{linear\_solve\_cholesky}(A,b_i)$ is
computing the Cholesky decomposition of the matrix $A$, which dominates the cost of the
subsequent forward substitution and backward substitution. Thus, an efficient way of
solving the linear systems (6.50) is to compute the Cholesky factor $U$ of $A$ only
once, outside the ``for'' loop, and then do the forward and backward substitutions for
solving each linear system inside the ``for'' loop; see the pseudocode from Table 6.3
for details.

Recall that both the forward substitution $\text{forward\_subst}(U\tp,b_i)$ and the
backward substitution $\text{backward\_subst}(U,y)$ require $n^2 + O(n)$ operations;
cf. (2.8) and (2.19). Thus, the operation count for solving the $p$ linear systems
using the method from Table 6.3 is
\[
\frac{1}{3}n^3 + O(n^2) \;+\; p\left(n^2+O(n) + n^2+O(n)\right)
\;=\; \frac{1}{3}n^3 + 2pn^2 + O(n^2) + pO(n),
\]
which is smaller than $\frac{1}{3}pn^3 + pO(n^2)$, the operation count required by
solving the $p$ linear systems sequentially; see (6.51).

As an example of solving multiple systems corresponding to the same matrix, we show
how to efficiently compute the inverse of a symmetric positive definite matrix.

\begin{table}[H]
\centering
\caption{Solution of multiple linear systems corresponding to the same spd matrix}
\fbox{\begin{minipage}{0.85\linewidth}
Input:\\
$A = $ symmetric positive definite matrix of size $n$\\
$b_i = $ column vectors of size $n$, $i = 1:p$
\bigskip

Output:\\
$x_i = $ solution to $Ax_i = b_i$, $i = 1:p$
\bigskip

\noindent
$U = \text{cholesky}(A)$\\
for $i = 1:p$\\
\hspace*{1.5em}$y = \text{forward\_subst}(U\tp,b_i)$;\\
\hspace*{1.5em}$x_i = \text{backward\_subst}(U,y)$;\\
end
\end{minipage}}
\end{table}

Let $A$ be a symmetric positive definite matrix of size $n$, and let
$A^{-1} = \text{col}(c_k)_{k=1:n}$ be the column form of the inverse matrix of $A$.
Then, $AA^{-1} = I$, where $I = \text{col}(e_k)_{k=1:n}$ is the identity matrix of size
$n$. Using (1.11), we find that $AA^{-1} = \text{col}(Ac_k)_{k=1:n}$, and therefore
$AA^{-1} = I$ is equivalent to
\[
Ac_k \;=\; e_k, \quad \forall\, k = 1:n,
\]
which can be solved using the method from Table 6.3 as follows:
\begin{center}
\fbox{\begin{minipage}{0.55\linewidth}
$U = \text{cholesky}(A)$\\
for $k = 1:n$\\
\hspace*{1.5em}$y = \text{forward\_subst}(U\tp,e_k)$\\
\hspace*{1.5em}$c_k = \text{backward\_subst}(U,y)$\\
end\\
$A^{-1} = \text{col}(c_k)_{k=1:n}$
\end{minipage}}
\end{center}

The operation count for computing the inverse matrix using this method is
\begin{equation}
\frac{1}{3}n^3 + O(n^2) \;+\; n\left(n^2+O(n) + n^2+O(n)\right) \;=\; \frac{7}{3}n^3 + O(n^2);
\end{equation}
see (10.89) from Section 10.2.3 for a proof of the last equality.


\section{Optimal linear solvers for tridiagonal symmetric positive definite matrices}

Solving linear systems corresponding to tridiagonal symmetric positive definite
matrices is often required in financial applications, e.g., for the efficient
implementation of the cubic spline interpolation and for the finite difference
solution of the Black--Scholes PDE; see section 6.4 and the references from section
6.5. Note that a linear system $Ax = b$ corresponding to a tridiagonal symmetric
positive definite matrix $A$ can be solved by using a linear solver based on the
Cholesky decomposition of $A$, or by using a linear solver based on the LU
decomposition of the matrix $A$.
% This file ends here (end of PDF page 192 / folio 176), mid-section 6.3,
% right before the section continues onto PDF page 193 (out of range for
% this agent). No environment is left open at this point.
